% SPDX-License-Identifier: Apache-2.0
% covers: I2cInit
\datasheet{I2cInit}{The I2C master that resets and configures an HDMI encoder chip.}

\dsfacts{\code{lib/parts/src/hdmi.rs}}{\code{txhdl\_parts::hdmi::I2cInit}}%
{\code{//docs:hdmi}}

\subsection*{Function}

\code{I2cInit} holds the chip's reset low on \code{nreset}, releases
it, waits, and then writes a table of register values as I2C
transactions. The table is the SiI9134's configuration,
\code{SII9134\_WRITES}, with two entries. The unit drives the two
open-drain lines through \code{scl\_low} and \code{sda\_low}, and reads
the data line back on \code{sda\_in}. \code{done} rises when the table
is written, and \code{failed} is high when the chip did not acknowledge
a byte.

The unit has no bus. It writes this table, one transaction per entry.
\code{//docs:hdmi} and the source both state that no document the part
was written from gives the chip's registers. The values are
unconfirmed until checked against the SiI9134's documentation or the
board vendor's example.

\begin{center}\footnotesize
\begin{tabular}{@{}llllp{0.4\textwidth}@{}}
\toprule
Entry & Address & Register & Value & Meaning, unconfirmed \\
\midrule
0 & \code{0x72} & \code{0x08} & \code{0x35} & System control: out of
power down, with the input bus as the board wires it. \\
1 & \code{0x7a} & \code{0x2f} & \code{0x00} & The output as DVI: no
HDMI packets, only the picture. \\
\bottomrule
\end{tabular}
\end{center}

The address is the chip's I2C address as eight bits, write bit
included.

\subsection*{Parameters}

\begin{center}\footnotesize
\begin{tabular}{@{}lp{0.7\textwidth}@{}}
\toprule
Parameter & Meaning \\
\midrule
\code{DIV} & Cycles in a quarter of an I2C bit. A bit is four
quarters. \\
\code{HOLD} & Cycles the reset is held low, and cycles the unit waits
after releasing it. \\
\bottomrule
\end{tabular}
\end{center}

The tables below are for \code{I2cInit<63, 2520000>}, the master of
\code{hdmi/src/netlist.rs}. Its doc comment states that a quarter of 63
cycles at 25.2~MHz makes a bit at 100~kHz, and that 2\,520\,000 cycles
are a tenth of a second.

\dstables{I2cInit}

\subsection*{Behaviour}

The unit is two processes. The quarter clock counts \code{tick} from 0
to \code{DIV} less one and round again, every cycle, and drives
\code{done} from \code{written} and \code{failed} from \code{nak}. The
other is the whole job written as the sequence it is: the reset held,
the reset released, then a transaction per entry, a quarter of a bit
per wait. The lowering numbers its waits into a state register the unit
does not declare, \code{at\_wait}, keeps its loops' counts in
\code{for0} to \code{for5}, and holds the three lines it drives
between states in \code{nreset\_held}, \code{scl\_low\_held} and
\code{sda\_low\_held}; all are in the tables above. \code{byte} is
\code{byte\_rw} in the netlist, since SystemVerilog reserves the word.

\begin{itemize}
\item \code{nreset} is low for \code{HOLD} cycles, then high, and the
unit waits \code{HOLD} cycles more before the first transaction.
\item A transaction is the start, the address, the register and the
value, each a byte followed by its acknowledge, and the stop, every
step a whole number of quarters of \code{DIV} cycles.
\item At the start the data line falls while the clock is high, then
the clock falls. At the stop the clock rises while the data line is
low, then the data line rises.
\item In a data bit the data line is set while the clock is low, and
the clock is high through the two middle quarters, so the data changes
only while the clock is low. Bytes go out most significant bit first,
from \code{shift}.
\item In an acknowledge the unit releases the data line, and reads
\code{sda\_in} at the end of the clock's high half, as it pulls the
clock low again. A high line sets \code{nak}, which \code{failed}
follows. Nothing clears \code{nak}.
\item \code{scl\_low} or \code{sda\_low} high pulls its line low.
\item After the stop of the last entry, \code{written} is set,
\code{done} rises a cycle later, and the sequence waits for ever.
\end{itemize}

\subsection*{Verification}

\begin{itemize}
\item \code{the\_chip\_is\_configured\_over\_i2c}, in
\code{//lib/parts:parts\_test}, runs \code{I2cInit<3, 5>} against
\code{I2cDevice}, a model of a device that acknowledges every byte and
records what it was sent. It checks that \code{done} rises, that
\code{failed} stays low, that the reset is released at cycle 5, and
that the device received \code{SII9134\_WRITES}.
\item \code{ex\_hdmi} runs \code{I2cInit<4, 10>} against the same model
and asserts the same table, \code{done}, and \code{failed} low.
\item The build simulates that unit's netlist against the run under
nvc and Verilator: \code{//docs:hdmi\_sim\_hdmi\_i2c\_hdmi\_i2c\_tb\_test}
and \code{//docs:hdmi\_vsim\_hdmi\_i2c\_test}.
\end{itemize}

\subsection*{Used by}

The example \code{ex\_hdmi}. The demonstration for the AX7A200B:
\code{//hdmi:netlist} writes the Verilog of \code{hdmi\_i2c}, and
\code{hdmi/board/hdmi\_demo.v} instantiates it with open-drain pads on
the chip's I2C lines and \code{done} and \code{failed} on LEDs.
\code{//hdmi:demo\_synth} and \code{//hdmi:demo\_pnr} put the design
through Vivado to a bitstream. \code{//docs:hdmi} states that it has
not yet run on the board.

\subsection*{Limits}

There is one table, and its values are unconfirmed. The unit keeps
writing after a byte is not acknowledged, and it writes the table once.
It has no input for the clock line. \code{tick} is 16 bits wide, so
\code{DIV} is at most 65\,536; the counters of the two holds are as
wide as \code{HOLD} needs, 22 bits in the tables.
